\section{Introduction}

\subsection{Project Overview}
Football coaches need to study two sides of a training session. Video shows
where players and the ball moved. Wearable sensors show the physical effort
behind that movement. When both records use different clocks and player names,
it is difficult to connect a hard run or rise in pulse with the event on the
field.

This project develops \systemname. A fixed camera records a controlled football
area. The vision pipeline detects players, goalkeepers, referees, and the ball.
ByteTrack keeps short-term track IDs, while team colour, appearance, jersey
votes, and manual review help resolve a track to a roster player.

Each selected player can wear an ESP32-S3 unit with an MPU-6050 inertial sensor,
MAX30102 optical sensor, and NEO-6M GPS receiver. The tested
\texttt{wearable\_stream} firmware is used as the main field build. It sends a
device ID, sequence number, source time, raw sensor samples, and status flags
to the required VPS relay. The local processor reads the relay stream and adds
quality information before any value is shown to the coach.

The two streams are joined by time and player identity. SNTP gives the wearable
a UTC time anchor, and its monotonic timer keeps sample order between clock
updates. Player-device binding states which
wearable belongs to which visible player. A joined record is created only when
the time error, packet age, signal quality, and identity confidence pass their
limits. A failed check gives a clear missing or low-confidence state instead of
a normal-looking value.

The project PC runs video processing, sensor processing, fusion, analytics,
local session files, the API, and the dashboard. Live network access is handled
through the required relay at \texttt{cpg44.nivaspms.com}. Cloudflare resolves
the domain and Caddy routes it to the CPG44 relay on the VPS. The relay keeps a
small memory-only replay window for short reconnects and stores no match
history.

The dashboard is arranged around a coach's work. The live page gives most of
the screen to a readable pitch view. It shows player tracks, identity state,
team, jersey label, ball path, possession evidence, and marked tactical areas.
The selected-player panel shows speed, distance, acceleration load, pulse
trend, estimated SpO$_2$, GPS state, packet age, and signal quality.

The analysis page links every alert to the same session timeline. Team width,
length, centroid, compactness, line spacing, pressure events, and high-intensity
efforts are calculated from pitch positions. Workload change is compared with
the player's earlier sessions when enough valid data exists. This lets a coach
inspect the reason behind a result instead of seeing only a score.

The platform separates measurement and interpretation. Detection, tracking,
pitch mapping, and sensor processing first create checked observations. Simple
formulas create movement, workload, and team-shape features. Rules and trained
models use these features only after their inputs pass the data-quality check.
Injury-related output is a workload warning for review, not a diagnosis
\cite{leckey2025injury}.

Public datasets support early training. SoccerNet-Tracking provides football
objects, team information, jersey metadata, and difficult occlusion cases
\cite{cioppa2022soccernet}. SportsMOT gives further examples of fast movement
and similar-looking players \cite{cui2023sportsmot}. Campus footage is still
needed because its camera height, field markings, kits, and lighting are
different.

At the mid-semester stage, the repository contains the vision and wearable
paths, fusion services, training scripts, API, live dashboard, relay service,
and evaluation plans. The tested \texttt{wearable\_stream} firmware compiles in
the \texttt{soccer} conda environment. Current software tests and detector logs are reported as
implementation evidence. Final accuracy will be based on the controlled campus
tests described later in this report.

\subsection{Need Analysis}
The main users are a coach and a student analyst. They need quick answers to
practical questions: Which player reaches a press late? When does the team
become too wide? Does a fall in running output occur with a rise in physical
load? Is a result weak because the ball, player identity, or sensor signal was
uncertain?

A camera alone cannot measure the player's body response. A wearable alone
cannot show team spacing, opponent pressure, or the ball. The proposed system
brings both views to one timeline and keeps the source evidence. This makes the
result easier to check during a campus training session.

The work is also useful as an evaluation platform. It records clock offset,
sequence gaps, mapping error, track identity, signal quality, corrections, and
feature versions. The team can therefore locate an error and compare changes
instead of judging the system only from a live screen.

\subsection{Research Gaps}
The following gaps were found from the literature and the prototype review.

\begin{enumerate}
  \item \textbf{Tactical and physical data are separated.} Football tracking
  work usually studies trajectories, while workload work uses GPS or training
  records. A useful campus system must join both at player and event level
  \cite{rossi2018injury,honda2022pass}.

  \item \textbf{Small and blocked objects remain difficult.} Ball size, camera
  motion, occlusion, and similar kits continue to affect football vision
  systems \cite{zheng2025sportsvision}. Short recovery is useful, but the
  system should not invent a ball position or player name.

  \item \textbf{Track ID is not the same as player identity.} ByteTrack can
  recover low-score detections, but it does not identify a football player
  \cite{zhang2022bytetrack}. Team evidence, jersey votes, appearance, and
  analyst confirmation are still required.

  \item \textbf{Wearable values can look valid during poor contact.} PPG is
  affected by motion and sensor placement \cite{pollreisz2022wearable,
  zhang2021ppg}. Heart rate and SpO$_2$ therefore need motion-aware quality
  checks and comparison with a reference device.

  \item \textbf{Time alignment is often hidden.} Camera time, microcontroller
  time, GPS time, and network arrival time are not equal. A joined result needs
  a measured timing error and visible sequence gaps.

  \item \textbf{Injury models do not transfer directly.} Studies use different
  players, injury definitions, and prediction windows
  \cite{pilka2023workload,leckey2025injury}. A local model needs repeated
  sessions, later-session testing, and clear limits.

  \item \textbf{Many tactical results are hard to explain.} Tactical models
  can use player relations and event context \cite{tacticai2024,
  teixeira2025tactical}. Coaches still need the supporting field position,
  time, input quality, and simple reason.
\end{enumerate}

Our work builds on these studies by joining the visible action, wearable load,
identity evidence, and time error in one record. It differs by focusing on a
complete and testable campus workflow, from data capture to live coach review.

\subsection{Problem Definition and Scope}
The problem is to build and evaluate a real-time football platform that tracks
players and the ball, reads wearable sensors, maps all values to one session
clock, and gives useful tactical and workload information. The system must keep
uncertain identity, timing, or sensor values visible.

The scope covers controlled football recordings, one calibrated camera,
ESP32-S3 wearables, local processing, required remote relay transport, live and
review dashboards, data export, and model training. Medical diagnosis, official
referee decisions, and clinical SpO$_2$ use are outside the scope.

\subsection{Assumptions and Constraints}
The camera is assumed to stay fixed after calibration and to show the marked
test area. Players use known team kits and a prepared roster. Wearables are
paired before the session and fitted safely.

The project uses one available PC, low-cost sensor modules, campus footage, and
a limited number of repeated sessions. Ball size, player overlap, network
jitter, PPG motion noise, and GPS dropout are expected constraints. The relay
must fit the VPS setup, use the domain name, use only a few megabytes of cache,
and create no persistent storage.

\subsection{Standards}
FIFA's Electronic Performance and Tracking Systems guidance separates product
safety from performance testing \cite{fifaepts}. We follow the same rule. A
wearable that powers on safely is not called accurate until it is compared with
reference measurements.

Manufacturer documents guide the sensor setup. The MAX30102 document gives its
sampling and optical settings \cite{max30102}. The MPU-6050 document gives IMU
ranges, filtering, and interfaces \cite{mpu6050}. The NEO-6 data sheet gives GPS
messages and update limits \cite{neo6}. IEEE numeric style is used for
references.

The live relay uses HTTPS and WebSocket transport, a private token, bounded
message size, value checks, and a fixed memory cache. No password or relay token
is kept in tracked firmware source or a frontend bundle.

\subsection{Approved Objectives}
The objectives approved in the proposal are:

\begin{enumerate}
  \item To develop an integrated data-processing framework that combines
  vision-based player tracking with wearable sensor streams for football
  monitoring.
  \item To use wearable hardware to collect physiological and spatial data
  from GPS, accelerometer, pulse-oximeter, and related sensors.
  \item To design machine-learning models for player and team performance
  trends, and a tactical module that gives recommendations from the analysed
  performance.
  \item To evaluate the system on available datasets and controlled test
  scenarios.
\end{enumerate}

\subsection{Methodology}
The work follows a measured pipeline. First, public and campus footage is
labelled and split by sequence. The detector, ByteTrack association, player
identity, ball recovery, and pitch mapping are then tested separately.

In parallel, the wearable sends timestamped IMU, GPS, and raw PPG samples
through the relay. The project PC checks these samples and estimates BPM and
SpO$_2$ only from acceptable PPG windows. Time alignment and player-device
binding create a joined event stream. The dashboard uses this stream for live
tracking and simple coach-facing analytics. Checked exports are later used for
model training and evaluation on unseen sessions.

\subsection{Project Outcomes and Deliverables}
The planned product is a working local football intelligence platform with a
required memory-only VPS relay. Deliverables include the wearable prototype,
firmware, vision and tracking pipeline, identity and time-sync services, live
dashboard, tactical and workload features, training pipeline, test results,
technical report, and demonstration.

\subsection{Novelty of Work}
The novelty is in the joined workflow, not in claiming a new sensor or tracker.
One event can carry player position, team context, ball state, physical load,
signal quality, identity confidence, and timing error. Coaches can correct an
identity and see why a tactical or workload warning was raised.

The platform also joins local processing with a storage-free live relay. The
same timestamped records support the dashboard, field validation, and later
model training. This gives a practical link between a live coaching tool and a
repeatable research dataset.
